% HAND-WRITTEN fragment -- not generated by maestro2tex, but placed by it via the
% `order` list in configs/tb_anatop_pixel_dacR2R12_corners.json.
%
% Numbers here come from tab_dacr2r12_reference, tab_dacr2r12_psu,
% tab_dacr2r12_mc_reference and tab_dacr2r12_mc_psu. Re-read them from those tables
% if either run is repeated.
%
% The ideal-ladder reference-current model and its fit were derived from the netlist
% topology and the nominal tight-tolerance sweep
% (simarchive/.../20260816_setup_and_verify/sweep_nominal_tight_tol.txt):
% I_ref maximum 1.4444 Vref/R at code 1365, r = -0.9698, R^2 = 94.0 %,
% residual 931 uV rms against a 25.700 mV spread, V0 = 2.1679 V.

\subsubsection{Reference movement and the supply-block requirement}
\label{sss:dacr2r12-supply}

Divide the reference out and the ladder's own transfer is well behaved; leave it in and
the same sweep is \SI{-37}{LSB} of INL and \SI{-10.8}{LSB} of DNL
(Table~\ref{tab:dacr2r12-wpsu}). The difference is entirely reference movement, and the
mechanism is not the one a reader expects.

In an R-2R the binary weighting comes from ladder attenuation, not from branch current:
every \(2R\) branch switched to the reference draws a comparable current. The current
drawn from \texttt{VddDac} therefore tracks the bit \emph{pattern}, not the code value.
An ideal-ladder model gives a reference current of zero at code~0,
\(0.25\,V_{\mathrm{ref}}/R\) at code~2048 and \(0.333\,V_{\mathrm{ref}}/R\) at code~4095,
with a maximum of \(1.444\,V_{\mathrm{ref}}/R\) at code~\(1365 =
\mathtt{0b010101010101}\) --- the alternating-bit pattern, and exactly where the measured
\texttt{VddDac} minimum sits, in every corner but ff and in all 200 Monte Carlo samples.
A one-parameter fit \(V_{\mathrm{DDDAC}} = V_0 - R_{\mathrm{out}} I_{\mathrm{ref}}\)
accounts for \SI{94}{\percent} of the variance of the measured curve, at a correlation of
\num{-0.97} and a residual of \SI{931}{\micro\volt} rms against a \SI{25.70}{\milli\volt}
spread. Figure~\ref{fig:dacr2r12-vddac} is that curve. It closes on the measured
nonlinearity: droop times code weight accounts for \SI{-36.9}{LSB} of the
\SI{-37.1}{LSB} of raw INL, so \SI{99.5}{\percent} of the delivered block's static error
is reference modulation and not a ladder error.

Two consequences follow that are easy to get wrong. First, the
\SI{-10.86}{LSB} DNL is real non-monotonicity of the block as delivered and not a
numerical artifact: tightening \texttt{reltol} from \num{1e-3} to \num{1e-6},
\texttt{vabstol} from \num{1e-6} to \num{1e-9} and \texttt{iabstol} from \num{1e-12} to
\num{1e-15} moves \texttt{VddDac} by at most \SI{193}{\micro\volt} anywhere and leaves the
structure of the curve intact. Second, endpoint-anchored INL makes the block look worse
still, because codes 0 and 4095 both sit near the top of the droop curve and the whole
bow therefore lands in INL at maximum leverage.

The requirement on the supply block follows from the DAC rather than from an assumed
target. A quarter-LSB reference budget is \(0.25 \times \SI{528.96}{\micro\volt} =
\SI{133}{\micro\volt}\), against the \SI{25.8}{\milli\volt} the ladder's own
code-dependent current actually produces: the supply's output impedance has to fall by a
factor of \num{193}, to about \SI{5}{\ohm}. Measured over the 200 Monte Carlo samples it
is \SI{894}{\ohm} small-signal at zero load (Table~\ref{tab:dacr2r12-mc-psu-zout}) and
\SI{1133}{\ohm} as a chord across the full \SIrange{0}{100}{\micro\ampere} load sweep
(Table~\ref{tab:dacr2r12-mc-psu-load}); the typical corner gives \SI{902}{\ohm} and
\SI{1137}{\ohm}. That is what the topology of
Figure~\ref{fig:dacr2r12-psu} gives: the block is a shunt clamp, not a regulated loop, so
its output impedance is the \(1/g_m\) of the shunt device at a quiescent current of
\SI{160}{\micro\ampere} and is flat from DC to \SI{100}{\kilo\hertz}. Closing the factor
of 193 is a topology change, not a trim. \texttt{VddDac} \emph{is} the ladder
reference today, not merely a gated supply for the digital side: \texttt{Out}/\texttt{VddDac}
is \num{0.499890} at code 2048 and \num{0.999733} at code 4095, against an ideal
\(4095/4096 = \num{0.999756}\).

The supply-side entries that follow are one defect measured five ways, not five problems:
\texttt{vddac\_pp}, the raw \texttt{dnl\_min}, \texttt{load\_reg}, \texttt{dvddac\_load}
and \texttt{zout\_max} are all the same output impedance. \texttt{load\_reg} and
\texttt{zout\_lf} correlate at \(r = \num{0.97}\) across the 200 samples, and
\texttt{vddac\_pp} against the raw \texttt{dnl\_min} at \(r = \num{-0.94}\).
